\documentclass[10pt,a4paper]{amsart}
\usepackage[margin=19mm]{geometry}
\usepackage{amsmath,amssymb,mathtools}
\usepackage[scr=rsfs,cal=euler]{mathalfa}
\usepackage{xfrac}
\usepackage[unicode,hidelinks,bookmarks=false]{hyperref}
\newcommand{\ie}{i.e.\ }
\newcommand{\cf}{cf.\ }
\newcommand{\resp}{respectively\ }
\newcommand{\Subsection}{\subsection}
\providecommand{\nobreakdash}{\nobreak\hbox{-}}
\setlength{\emergencystretch}{2em}
\multlinegap0pt
\usepackage{bm}
\usepackage{bbm}
\usepackage{dsfont}
\usepackage{xspace}
\usepackage{cancel}
\usepackage{mathtools}
\usepackage{amsthm}
\makeatletter
\@ifundefined{Proposition}{\newtheorem{Proposition}{Proposition}}{}
\@ifundefined{lemma}{\newtheorem{lemma}{Lemma}}{}
\makeatother
\title[Soliton resolution, asymptotic stability and Painlev\'e tran]{Soliton resolution, asymptotic stability and Painlev\'e transcendents in the combined Wadati-Konno-Ichikawa and short-pulse equation}
\author{Yidan Zhang, Engui Fan}
\date{Source: arXiv:2505.03144v3 (CC BY 4.0)}
\begin{document}
\maketitle

\noindent\emph{Source and licence.} Soliton resolution, asymptotic stability and Painlev\'e transcendents in the combined Wadati-Konno-Ichikawa and short-pulse equation. Version arXiv:2505.03144v3 (CC BY 4.0), distributed under the Creative Commons Attribution 4.0 International licence, as recorded in the arXiv licence metadata for this exact version and shown on the abstract page https://arxiv.org/abs/2505.03144v3. Full rights notice: licenses/arxiv-2505.03144-CC-BY-4.0.txt.\par\smallskip

\noindent\textit{Editorial scope.} Counted unit: the Proposition whose source label is \texttt{p17} in the article (source0.tex lines 2489--2495, source label \texttt{p17}), delivered together with the article's own complete original proof of it (source0.tex lines 2497--2570, one \texttt{proof} environment of 74 lines). No mathematics is written, repaired or re-derived: statement and proof are the article's own text line for line. Required context retained uncounted: Imthata (lines 847--849); p10 (lines 1373--1379); p11 (lines 1384--1405); eq36 (lines 2480--2482). Every definition, display and label of those lines is retained unchanged. Omitted from this package and not redistributed: 1--846, 850--1372, 1380--1383, 1406--2479, 2483--2488, 2496--2496, 2571--4695. Nothing inside a retained excerpt is omitted or altered: every retained body is delivered exactly as the article's lines are written, so no omission receipt of any kind accompanies this package. Citations: the retained excerpts carry no citation command at all, so no bibliography entry of the article is transcribed and no reference is redistributed. Reference adaptation: none was needed and none was made; every reference inside the delivered unit resolves inside it, so no unresolved reference or citation is produced. Printed slips kept literal: no printed word, symbol, hypothesis or number of the counted statement or of its proof was altered. Layout and class: the article's own class is \texttt{article} with packages including color, graphicx, amsmath, appendix, float, amssymb, enumitem, mathrsfs, natbib, amsthm, pgf, CJK, tikz, stmaryrd; the wrapper is the lane's pinned \texttt{amsart} editorial wrapper at 10pt on A4, which restates the theorem-like environments the retained text needs and adds the article's own macro definitions that the retained text needs (none), with extra wrapper packages (bm, bbm, dsfont, xspace, cancel, mathtools). The delivered numbering is the unit's own. Assets: no retained span contains a figure environment, a graphics-inclusion command, a PDF-page inclusion, a table or a TikZ picture; no image file is copied into this package, the package declares no asset dependency and no image file is redistributed with it. Licence: version arXiv:2505.03144v3 of this article is distributed under the Creative Commons Attribution 4.0 International licence, as recorded in the arXiv licence metadata for this exact version and shown on the abstract page https://arxiv.org/abs/2505.03144v3. A full rights notice is delivered at licenses/arxiv-2505.03144-CC-BY-4.0.txt. Characters: every retained excerpt is pure ASCII, so the delivered engine, XeLaTeX with its default Latin Modern text font, reports no missing character.
\begin{align}
    \theta^{'}(k)&=\xi+12\alpha k^2+\frac{\beta}{4k^2},\nonumber\\ \mathrm{Im}\theta(k)&=\mathrm{Im}k\left[\xi+12\alpha|k|^2-16\alpha (\mathrm{Im}k)^2+\frac{\beta}{4|k|^2} \right].\label{Imthata}
\end{align}

\begin{lemma}\label{p10}\textup{(}near saddle points $k_j$\textup{)}
     $\mathrm{Im}\theta(k)$ defined by \eqref{Imthata} has the following estimates:
    \begin{align*}
        \mathrm{Im}\theta(k)&\gtrsim |\mathrm{Im}(k)|\left|\mathrm{Re}k-k_j\right|,\quad k\in\Omega_1,\quad j=1,\dots,\Lambda,\\
\mathrm{Im}\theta(k)&\lesssim -|\mathrm{Im}(k)|\left|\mathrm{Re}k-k_j\right|,\quad k\in\Omega_2,\quad j=1,\dots,\Lambda.
    \end{align*}
\end{lemma}

\begin{Proposition}\label{p11}
There exist  the functions $R_{\ell}(k)$: $\bar{\Omega}_{\ell}\to \mathbb{C}$, $\ell=1,2$  with the boundary values
	\begin{align} &R_{1}(k)=\left\{\begin{array}{ll}
	\rho(k)T_+(k)^{-2}, &k\in \mathbb{R},\\[4pt]
 \rho(k_j)T_0(k_j)^{-2} (k-k_j)^{-2\eta(k_j)i\nu(k_j)},  &k\in \Sigma_{1},\\
	\end{array}\right.\label{eq4} \\[5pt]
	&R_{2}(k)=\left\{\begin{array}{ll} \bar{\rho}(k)T_-(k)^{2},
 &k\in  \mathbb{R},\\[4pt] \bar{\rho}(k_j)T_0(k_j)^{2}(k-k_j)^{2\eta(k_j)i\nu(k_j)}, &k\in \Sigma_{2},
	\end{array} \right.
	\end{align}	
where $j=1,\cdots, \Lambda$.
The functions  $R_{\ell}(k), \ell=1,2$ admit the following estimates:
	\begin{align}
    &|R_{\ell}(k)|\lesssim 1+\left[1+\mathrm{Re}^2(k)\right]^{-\frac{1}{2}},\quad\;for\ k \in\Omega,\label{eq5}\\
    &|\bar{\partial}R_{\ell}(k)|\lesssim\chi(\mathrm{Re}k)+|r^\prime(\mathrm{Re}k)|+|k-k_j|^{-\frac{1}{2}}, \quad   \;for\  k\in \Omega,\;j=2,3 \;of\;case\;\textup{\uppercase\expandafter{\romannumeral1}},\label{eq6}\\
     &|\bar{\partial}R_{\ell}(k)|\lesssim \chi(\mathrm{Re}k)+|r'(\mathrm{Re}k)|+|k-k_j|^{-\frac{1}{2}}, \quad \; for\   k\in \Omega, \;j=1,2\;of\;case\;\textup{\uppercase\expandafter{\romannumeral4}},\label{eq7}\\
	&|\bar{\partial}R_{\ell}(k)|\lesssim|r'(\mathrm{Re}k)|+|k-k_j|^{-\frac{1}{2}}, \quad  for\   k\in \Omega, \;j=1,4\;of\;case\;\textup{\uppercase\expandafter{\romannumeral1}},\label{eq8}\\
&|\bar{\partial}R_{\ell}(k)|\lesssim|k| \quad as\ k\to0,\;for\ k\in \Omega,\label{eq9}\\
&\bar{\partial}R_{\ell}(k)=0,\quad \;for\ k\in\mathbb{C}\setminus\Omega,\nonumber
	\end{align}
where $\chi \in C^{\infty}_0(\mathbb{R},[0,1])$ is a fixed cut-off function with support near 0.
\end{Proposition}

\begin{equation}\label{eq36}
    S\left[f\right](k) = -\frac{1}{\pi} \iint_\mathbb{C} \frac{f(s) W^{(3)}(s)}{s - k} \, \mathrm{d}A(s),
\end{equation}

\begin{Proposition}\label{p17}
    Consider the operator $S$ defined by \eqref{eq36},  we can obtain $S : L^\infty(\mathbb{C}) \to L^\infty(\mathbb{C}) \cap C^0(\mathbb{C})$ and
\begin{equation}
    \|S\|_{L^\infty(\mathbb{C}) \to L^\infty(\mathbb{C})} \lesssim t^{-\frac{1}{4}}.
\end{equation}

\end{Proposition}

\begin{proof}
    For any $f \in L^\infty$, we have
\begin{equation*}
    \|Sf\|_{L^\infty} \leq \|f\|_{L^\infty} \frac{1}{\pi} \iint_\mathbb{C} \frac{|W^{(3)}(s)|}{|s - k|} \, \mathrm{d}A(s).
\end{equation*}
Recalling our definition $W^{(3)} =  M_{RHP}^{(2)}(k) \bar{\partial} R^{(2)}(k) M_{RHP}^{(2)}(k)^{-1}.$ First we know that $W^{(3)}(k) \equiv 0$ for $k \in \mathbb{C} \setminus \bar{\Omega}$. Besides, we only take into account the matrix-valued functions have support in sector $\bar{\Omega}$. Moreover, we know that $M_{RHP}^{(2)}(k)$ and $M_{RHP}^{(2)}(k)^{-1}$ are all bounded on $\bar{\Omega}$, which means
\begin{equation}\label{eq37}
    \iint_{\Omega_\ell}\frac{|W^{(3)}(s)|}{|s-k|}\mathrm{d}A(s)\lesssim\iint_{\Omega_\ell}\frac{|\bar{\partial}R_\ell(s)e^{\pm2it\theta}|}{|s-k|}\mathrm{d}A(s),\quad\ell=1,2,
\end{equation}
where the superscript takes $+$ for $\ell=1$, takes $-$ for $\ell=2$.
To shorten the length of this paper, we only consider the region $\Omega_1\cap\left\{k\in\mathbb{C}:\mathrm{Re}k>k_1\right\}:=\widehat\Omega_1$ of case \uppercase\expandafter{\romannumeral1}. Together with Proposition \ref{p11}, we can break right side of the equation \eqref{eq37} into two parts:
\begin{equation*}
\iint_{\widehat\Omega_1} \frac{|\bar{\partial} R_{1}(s)| e^{-2t \mathrm{Im} \theta}}{|k - s|} \, \mathrm{d}A(s)
\lesssim L_1 + L_2,
\end{equation*}
with
\begin{align*}
    L_1 = \iint_{\widehat\Omega_1} \frac{|r'(\mathrm{Re}s)| e^{-2t \mathrm{Im} \theta}}{|k - s|} \, \mathrm{d}A(s),\quad L_2 = \iint_{\widehat\Omega_1} \frac{|s - k_1|^{-\frac{1}{2}} e^{-2t \mathrm{Im} \theta}}{|k - s|} \, \mathrm{d}A(s).
\end{align*}
Denote $k=x+yi,s=k_1+u+iv$ with $x,y,u,v\in\mathbb{R}$, then Lemma \ref{p10} implies that
\begin{align*}
L_1&\lesssim\int_0^{+\infty}\int_v^{+\infty}\frac{|r'(\mathrm{Re}s)| e^{- t uv}}{|k-s|}\mathrm{d}u\mathrm{d}v\leqslant\int_0^{+\infty}e^{- t v^2}\mathrm{d}v\int_v^{+\infty}\frac{|r^\prime(k_1+u)| }{|k-s|}\mathrm{d}u\\
    &\leqslant \int_0^{+\infty}e^{- t v^2}\|r^\prime\|_{L^2}\|\frac{1}{|k-s|}\|_{L^2(v,+\infty)}\mathrm{d}v\lesssim  \int_0^{+\infty}e^{- t v^2}\|\frac{1}{|k-s|}\|_{L^2(v,+\infty)}\mathrm{d}v.
\end{align*}
For further calculation, we introduce the following estimate for $q>1$,
\begin{equation}\label{eq41}
    \begin{aligned}
    \|\frac{1}{|k-s|}\|_{L^q(v,+\infty)}&=\left(\int_v^{+\infty}\frac{1}{|k-s|^q}\mathrm{d}u\right)^{\frac{1}{q}}\\
    &\leqslant|v-y|^{\frac{1}{q}-1}\int_0^{+\infty}\left[\left(\frac{u+k_1-x}{v-y}\right)^2+1\right]^{-\frac{q}{2}}\mathrm{d}\left(\frac{u+k_1-x}{v-y}\right)\\
    &\lesssim|v-y|^{\frac{1}{q}-1}.
\end{aligned}
\end{equation}


Then back to the calculation of $L_1$, we have
\begin{equation}
    L_1\lesssim\int_0^{+\infty}\frac{e^{-t v^2}}{\sqrt{|v-y|}}\mathrm{d}v=L_1^{(1)}+L_1^{(2)},
\end{equation}
where
\begin{equation*}
     L_1^{(1)}=\int_0^{y}\frac{e^{-t v^2}}{\sqrt{y-v}}\mathrm{d}v,\quad L_1^{(2)}=\int_y^{+\infty}\frac{e^{-t v^2}}{\sqrt{v-y}}\mathrm{d}v.
\end{equation*}
Therefore,
\begin{align*}
     L_1^{(1)}\lesssim t^{-\frac{1}{4}}\int_0^1\frac{\mathrm{d}m}{\sqrt{m(1-m)}}\lesssim t^{-\frac{1}{4}},\quad L_1^{(2)}\lesssim \int_0^{+\infty}\frac{e^{-tm^2}}{\sqrt{m}}\mathrm{d}m\lesssim t^{-\frac{1}{4}},
\end{align*}
which implies $L_1\lesssim t^{-\frac{1}{4}}$.\\
As for $L_2$, by H\"older inequality with $\frac{1}{p} + \frac{1}{q} = 1,p>2$,
\begin{align}\label{eq38}
    L_2&\lesssim \int_0^{+\infty}e^{-t v^2}\|\frac{1}{\sqrt{|s-k_1|}}\|_{L^p(\mathbb{R}_+)}\|\frac{1}{k-s}\|_{L^q(\mathbb{R}_+)}\mathrm{d}v,
\end{align}
where
\begin{align*}
    \|\frac{1}{\sqrt{|s-k_1|}}\|_{L^p(\mathbb{R}_+)}&=\left(\int_0^{+\infty}\left(u^2+v^2\right)^{-\frac{p}{4}}\mathrm{d}u\right)^{\frac{1}{p}}\\
    &=v^{\frac{1}{p}-\frac{1}{2}}\left[\int_0^{+\infty}\left(1+m^2\right)^{-\frac{p}{4}}\mathrm{d}m\right]^{\frac{1}{p}}\lesssim v^{\frac{1}{p}-\frac{1}{2}}.
\end{align*}
Taking this estimate into equation \eqref{eq38}, we obtain
\begin{equation*}
    L_2\lesssim\int_0^{+\infty}e^{-t v^2}v^{\frac{1}{p}-\frac{1}{2}}|v-y|^{\frac{1}{q}-1}\mathrm{d}v=L_2^{(1)}+L_2^{(2)},
\end{equation*}
where
\begin{equation*}
    L_2^{(1)}=\int_0^{y}e^{-t v^2}v^{\frac{1}{p}-\frac{1}{2}}(y-v)^{\frac{1}{q}-1}\mathrm{d}v,\quad L_2^{(2)}=\int_y^{+\infty}e^{-t v^2}v^{\frac{1}{p}-\frac{1}{2}}(v-y)^{\frac{1}{q}-1}\mathrm{d}v.
\end{equation*}
Let $v=my$,  $L_2^{(1)}$ becomes
\begin{align*}
    L_2^{(1)}=\int_0^1e^{- ty^2m^2}y^{\frac{1}{2}}m^{\frac{1}{p}-\frac{1}{2}}(1-m)^{\frac{1}{q}-1}\mathrm{d}m\lesssim t^{-\frac{1}{4}}\int_0^1m^{\frac{1}{p}-1}(1-m)^{\frac{1}{q}-1}\mathrm{d}m\overset{{p,q>2}}\lesssim t^{-\frac{1}{4}}.
\end{align*}
Let $n=v-y$,  $L_2^{(2)}$ becomes
\begin{align*}
    L_2^{(2)}=\int_0^{+\infty}e^{-t(y+n)^2}(y+n)^{\frac{1}{p}-\frac{1}{2}}n^{\frac{1}{q}-1}\mathrm{d}n\leqslant \int_0^{+\infty}\frac{e^{- tn^2}}{\sqrt{n}}\mathrm{d}n\lesssim t^{-\frac{1}{4}}.
\end{align*}
From the above calculation, we obtain $L_2\lesssim t^{-\frac{1}{4}}$. Summarizing the results we give above,  $\|S\|_{L^\infty(\mathbb{C}) \to L^\infty(\mathbb{C})} \lesssim t^{-\frac{1}{4}}$ as $t\to\infty$.
\end{proof}

\end{document}
